\documentclass[12pt]{article}

\usepackage[a4paper,margin=1in]{geometry}
\usepackage{amsmath,amssymb}
\usepackage{newtxtext,newtxmath} % Times-style font
\usepackage{fancyhdr}
\usepackage{listings}
\usepackage{graphicx}
\usepackage{float}
\usepackage{xcolor}
\usepackage{booktabs}
\usepackage{hyperref}
\usepackage{subcaption}

\setlength{\parindent}{0pt}
\setlength{\parskip}{5pt}

% Keep entire document at 12 pt
\AtBeginDocument{
    \fontsize{12pt}{14pt}\selectfont
}

% ------------------------------------------------
% Header & Footer
% ------------------------------------------------
\pagestyle{fancy}
\fancyhf{}
\fancyhead[L]{Name: \underline{Shaik Jnaved}}
\fancyhead[R]{Enrollment No.: \underline{P26DS001}}
\fancyfoot[C]{\thepage}
\renewcommand{\headrulewidth}{0.4pt}
\setlength{\headheight}{15pt}

% ------------------------------------------------
% Code formatting
% ------------------------------------------------
\lstset{
    language=Python,
    basicstyle=\ttfamily\fontsize{8pt}{10pt}\selectfont,
    keywordstyle=\color{blue}\bfseries,
    stringstyle=\color{teal},
    commentstyle=\color{gray}\itshape,
    numbers=left,
    numberstyle=\tiny\color{gray},
    stepnumber=1,
    numbersep=6pt,
    breaklines=true,
    frame=single,
    columns=fullflexible,
    keepspaces=true,
    showstringspaces=false,
    tabsize=4
}

\begin{document}

% ------------------------------------------------
% TITLE
% ------------------------------------------------
\begin{center}
\textbf{Computer Science and Engineering Department, SVNIT, Surat}\\
\textbf{M.Tech. I -- Semester I}\\
\textbf{Computer Vision and Image Processing (CS6111 / CSDS119)}\\
\textbf{Lab Assignment: Haar Cascade Face Detection, Model Explanation, and Empirical Perturbation Analysis}
\end{center}

\vspace{0.2cm}

\textbf{Problem Statement:} \\
For the assigned task, the student is expected to:
\begin{enumerate}
    \item Study and execute the given GitHub Haar Cascade face-detection implementation.
    \item Use an LLM for code-understanding, debugging, and explaining the role of the cascade model and detection parameters.
    \item Make the required modifications to create a modular, robust, and telemetry-enabled implementation.
    \item Empirically test the detector across at least three parameter settings under systematic variations in \textbf{scale}, \textbf{illumination}, and \textbf{face orientation}, quantitatively reporting successful and failed detections.
    \item Submit the complete code, visual outputs, LLM interaction transcript, and experimental observations as a formal report.
\end{enumerate}

\vspace{0.2cm}
\hrule
\vspace{0.2cm}

\section{Theoretical Foundation of Haar Cascade Face Detection}

The Viola-Jones face detection framework (2001) represents a seminal milestone in computer vision. It is constructed upon four fundamental concepts:

\begin{enumerate}
    \item \textbf{Haar-like Rectangle Features:} Rather than working directly with raw pixel intensities, the detector computes differences between adjacent rectangular regions (two-rectangle horizontal/vertical, three-rectangle, and four-rectangle diagonal filters). These capture anthropometric invariants of human faces (e.g., eye region darker than forehead; nasal bridge brighter than bilateral cheeks).
    
    \item \textbf{Integral Image Representation (Summed Area Table):} Given an image $I(x,y)$, the integral image is defined as:
    \begin{equation}
        II(x,y) = \sum_{x' \le x, \, y' \le y} I(x',y')
    \end{equation}
    The sum of pixels within any arbitrary rectangular bounding box with corners $A(x_1, y_1), B(x_2, y_1), C(x_1, y_2), D(x_2, y_2)$ is computed in strictly \textbf{$O(1)$ constant time} via:
    \begin{equation}
        \text{Sum}(ABCD) = II(D) + II(A) - II(B) - II(C)
    \end{equation}
    This enables real-time evaluation regardless of the filter scale.
    
    \item \textbf{AdaBoost Feature Selection:} A $24 \times 24$ sub-window yields over 160,000 possible Haar features. AdaBoost selects an optimal small ensemble of weak classifiers (decision stumps) $h_j(x)$ and combines them into a strong classifier $H(x) = \text{sign}\left(\sum_{t=1}^T \alpha_t h_t(x)\right)$.
    
    \item \textbf{Attentional Cascade Architecture:} To achieve real-time throughput, classifiers are arranged in a degenerate decision tree of sequential stages. Because $>99.9\%$ of image sub-windows are non-faces, early stages with very few features ($1-10$) reject the vast majority of negative windows immediately. Only candidate windows passing all stages are classified as faces.
\end{enumerate}

\section{Mathematical Role of Detection Parameters in OpenCV}

The OpenCV detection method \texttt{cv2.CascadeClassifier.detectMultiScale()} exposes several critical hyperparameters:

\begin{itemize}
    \item \textbf{\texttt{scaleFactor}} (e.g., $1.05, 1.10, 1.30$): Governs the multi-scale image pyramid scaling step. At each level $k$, the image is scaled by $S_k = (\text{scaleFactor})^{-k}$. Smaller values ($1.03-1.05$) create dense pyramid steps, catching faces of arbitrary intermediate sizes at the expense of higher computation time. Larger values ($1.30$) accelerate execution but risk stepping over faces whose true scale falls between pyramid levels.
    
    \item \textbf{\texttt{minNeighbors}} (e.g., $2, 4, 8$): Specifies the clustering consensus threshold. Multiple overlapping candidate windows typically detect a true face. Detections are clustered using rectangle grouping; \texttt{minNeighbors} specifies the minimum count of constituent rectangles required to retain a cluster. Lower values increase recall but introduce false positives; higher values guarantee high precision but miss faint or occluded faces.
    
    \item \textbf{\texttt{minSize} / \texttt{maxSize}}: Lower and upper bounds on candidate face window size in pixels. Prunes out impossible scale ranges, preventing spurious noise detection in background textures.
\end{itemize}

\newpage

\section{Code Redesign and Implementation Modifications}

The original implementation (\texttt{main.py}) suffered from critical limitations: hardcoded webcam input (\texttt{cv2.VideoCapture(0)}), hardcoded parameters (\texttt{1.1, 4}), unhandled GUI crashes in headless/WSL environments, zero telemetry, and brittle file paths. 

The improved script (\texttt{main\_improved.py}) introduces a robust, production-ready architecture:

\begin{lstlisting}[caption={Key segments of the enhanced face detection engine (main\_improved.py)}]
def detect_faces(img, cascade, scaleFactor=1.1, minNeighbors=4, minSize=(30, 30), maxSize=None):
    gray = cv2.cvtColor(img, cv2.COLOR_BGR2GRAY)
    kwargs = {
        'scaleFactor': float(scaleFactor),
        'minNeighbors': int(minNeighbors),
        'minSize': minSize
    }
    if maxSize is not None:
        kwargs['maxSize'] = maxSize
        
    start_t = time.perf_counter()
    faces = cascade.detectMultiScale(gray, **kwargs)
    elapsed_ms = (time.perf_counter() - start_t) * 1000.0
    return faces, elapsed_ms

def annotate_frame(img, faces, elapsed_ms, scaleFactor, minNeighbors):
    annotated = img.copy()
    w_img = annotated.shape[1]
    for idx, (x, y, w, h) in enumerate(faces):
        cv2.rectangle(annotated, (x, y), (x + w, y + h), (0, 255, 128), 2)
        # Corner accents and dimension tags
        label = f"Face #{idx+1} ({w}x{h})"
        cv2.putText(annotated, label, (x, y - 5), cv2.FONT_HERSHEY_SIMPLEX, 0.45, (0, 255, 255), 1)
    # HUD Banner
    hud_text = f"Faces: {len(faces)} | Latency: {elapsed_ms:.1f}ms | sF: {scaleFactor} | mN: {minNeighbors}"
    cv2.putText(annotated, hud_text, (10, 20), cv2.FONT_HERSHEY_SIMPLEX, 0.5, (255, 255, 255), 1)
    return annotated
\end{lstlisting}

\section{Experimental Design and Parameter Configurations}

We defined four distinct parameter configurations representing different operational regimes:

\begin{table}[H]
\centering
\caption{Tested Detector Parameter Configurations}
\begin{tabular}{@{}lcccc@{}}
\toprule
\textbf{Configuration} & \textbf{scaleFactor} & \textbf{minNeighbors} & \textbf{minSize} & \textbf{Operational Regime} \\
\midrule
\textbf{Config 1 (Baseline)} & 1.10 & 4 & $(30, 30)$ & Standard balanced default \\
\textbf{Config 2 (High Sensitivity)} & 1.05 & 2 & $(20, 20)$ & Permissive / High recall \\
\textbf{Config 3 (High Precision)} & 1.30 & 8 & $(40, 40)$ & Strict / Fast / High precision \\
\textbf{Config 4 (Fine Strict)} & 1.03 & 6 & $(30, 30)$ & Dense pyramid / Strict consensus \\
\bottomrule
\end{tabular}
\end{table}

\newpage

\section{Experimental Results and Comparative Analysis}

\subsection{Experiment 1: Scale Variation}
Testing across five scales (Tiny: 0.18$\times$, Small: 0.30$\times$, Medium: 0.50$\times$, Standard: 1.00$\times$, Large: 1.60$\times$):

\begin{figure}[H]
    \centering
    \includegraphics[width=0.95\textwidth]{images/scale_experiment_grid_opt.jpg}
    \caption{Scale variation experiment across Configs 1--4. Yellow boxes represent Ground Truth; green boxes represent True Positive detections with $\text{IoU} \ge 0.35$.}
\end{figure}

\textbf{Observations:} Config 3 fails at Tiny scale ($\sim 30$ px face) because of its \texttt{minSize=(40,40)} constraint. Config 2 succeeds with IoU 0.879 due to its lower \texttt{minSize=(20,20)} and permissive consensus. Latency increases by $>10\times$ between Config 3 (5.9 ms) and Config 4 (63.5 ms).

\subsection{Experiment 2: Illumination Variation}
Testing across severe underexposure, moderate dim light, normal diffuse, moderate overexposure, specular burn/flash, and directional side shadow:

\begin{figure}[H]
    \centering
    \includegraphics[width=0.95\textwidth]{images/illumination_experiment_grid_opt.jpg}
    \caption{Illumination variation experiment across six illumination regimes.}
\end{figure}

\textbf{Observations:} Linear illumination shifts (dim lighting, moderate overexposure, directional shadow) are well tolerated because Haar rectangle differences cancel out uniform offsets. However, under non-linear clipping (Severe Underexposure and Specular Burn), Config 3 fails completely due to loss of facial edge contrast.

\newpage

\subsection{Experiment 3: Orientation Variation}
Testing in-plane rotations from $0^\circ$ to $180^\circ$:

\begin{figure}[H]
    \centering
    \includegraphics[width=0.95\textwidth]{images/orientation_experiment_grid_opt.jpg}
    \caption{Orientation variation experiment showing detection breakdown beyond $\pm 15^\circ$.}
\end{figure}

\textbf{Observations:} Haar rectangle features are rigidly axis-aligned. Beyond $\pm 15^\circ$, facial symmetry and horizontal eye-band contrast are destroyed. Consequently, \textbf{all four detector configurations fail completely at $+30^\circ$, $+45^\circ$, $+90^\circ$, and $180^\circ$}.

\subsection{Quantitative Telemetry and Latency Benchmarks}

\begin{figure}[H]
    \centering
    \includegraphics[width=0.95\textwidth]{images/latency_and_accuracy_summary_opt.jpg}
    \caption{Left: Average latency (ms) per configuration. Right: Outcome distribution across all 18 test scenarios.}
\end{figure}

\begin{table}[H]
\centering
\caption{Quantitative Performance Summary Across All 18 Test Scenarios}
\begin{tabular}{@{}lccccc@{}}
\toprule
\textbf{Configuration} & \textbf{Mean Latency} & \textbf{Success (Exact)} & \textbf{Partial (TP+FP)} & \textbf{Failed (Missed)} & \textbf{Success Rate} \\
\midrule
Config 1 (Baseline) & 20.4 ms & 12 / 18 & 0 / 18 & 6 / 18 & 66.7\% \\
Config 2 (High Sensitivity) & 39.5 ms & 11 / 18 & 3 / 18 & 4 / 18 & 61.1\% \\
Config 3 (High Precision) & \textbf{6.1 ms} & 10 / 18 & 0 / 18 & 8 / 18 & 55.6\% \\
Config 4 (Fine Strict) & 53.6 ms & 11 / 18 & 3 / 18 & 4 / 18 & 61.1\% \\
\bottomrule
\end{tabular}
\end{table}

\section{Conclusion}
This investigation confirmed the theoretical principles and practical limits of the Viola-Jones Haar Cascade detector. While computationally lightweight and effective for upright frontal faces under moderate lighting, the model is inherently susceptible to extreme illumination saturation and in-plane head rotations exceeding $\pm 15^\circ$. For unrestricted real-world deployment, modern deep convolutional neural networks (e.g., MTCNN, RetinaFace, or YOLO) provide the necessary rotational invariance and feature expressiveness.

\end{document}
